\chapter{全书计划与配套材料}

本版完成前言与第一、二章初稿。后续章节按研究表达任务推进，每章包括 AI 可以辅助什么、SciDraw AI 操作路径、人工核验、修改案例与交付物。

\begin{longtable}{p{0.08\linewidth}p{0.38\linewidth}p{0.43\linewidth}}
\toprule
\textbf{章} & \textbf{主题} & \textbf{主要交付物} \\
\midrule
\endhead
1 & 建立 AI 科研绘图工作流 & 任务书、路线与资料目录；已写初稿 \\
2 & 研究内容转为绘图任务 & 关系表、输入与验收条件；已写初稿 \\
3 & 构图与草图 & 两个构图候选、选择理由 \\
4 & 视觉语言 & 颜色、标签、字体、箭头规范 \\
5 & 数据可视化 & 数据、成图、修订及核对记录 \\
6 & 方法与流程 & 步骤表、流程图、条件核对 \\
7 & 机制与结构 & 关系表、初稿、纠错与图例 \\
8 & 实验图像呈现 & 原始材料、面板计划、处理说明 \\
9 & 组图与图注 & 多面板图、图注、处理记录 \\
10 & 图形摘要 & 摘要任务书、布局候选、定稿 \\
11 & 具体而有效的修改 & 问题单、逐版修改和复核 \\
12 & 汇报与海报 & 同一研究的不同观看场景版本 \\
13 & 检查与交付 & 最终图、图注、来源、AI 使用记录 \\
\bottomrule
\end{longtable}

\section{可复用的任务书}
\begin{Prompt}
【用途与读者】
【一句话目标】
【素材与来源】
【必须出现的内容】
【关系与顺序】
【证据状态：观测、解释、假设、待确认】
【表达方案】
【不可改变的内容】
【交付要求与查阅日期】
【验收办法】
\end{Prompt}

\section{模拟数据说明}
\texttt{data/paired-example.csv} 是本书人为编制的教学数据。八个虚构对象各有两次测量，单位为任意单位；没有缺失值或预设剔除条件。七个对象上升，SIM04 从 18 降至 17。前测均值为 13.125，后测均值为 16.5。这些描述不支持真实研究、因果或疗效结论。

\texttt{scripts/plot\_paired.py} 从原始 CSV 绘制图像并核对上述性质。中文版和英文版分别生成带相应语言标签的 PDF 与 PNG。脚本不补写、删除或修改观测值。

\section{图版与 AI 使用记录}
机制图来自 Codex 内置图像生成，完整输入位于 \path{provenance/mechanism-prompt.txt}，核验记录位于同目录的 \path{mechanism-review.json}。模型的具体名称未由工具返回，本书不自行推断。该图保留原始输出，没有后期视觉编辑。

实际生成输入把关系解释放在图注，而非图内中文标签中；这项改动保留了关系含义，并让中英文版本共用图版。核验检查了对象、方向、线型和额外内容；节点的轻微渐变作为风格偏差记录。教学模型的关系是假定条件，不具有生物学验证意义。

工作流程图为 Codex 编写的 TikZ 图，界面图为真实开发版截图。各图均有来源与用途说明，具体清单见仓库的 ASSETS.md。后续加入真实研究案例时，应补齐来源、必要授权、实际操作和修订记录。
